\section{Sign patterns of functions}\label{app:modes}

We prove the lemmas used for Theorem~\ref{thm:two-modes}, with the notation
of Sections~\ref{sec:convex} and~\ref{sec:modes}. Deleting entries from a list
replaces its pattern by a subsequence. An alternating list $p$ satisfies
$p\preceq q$ for every alternating list $q$ with the same first sign and at
least as many entries.

\begin{lemma}[Descartes' rule for exponential sums]\label{lem:descartes}
Let $\lambda_1<\dots<\lambda_n$ and $b_1,\dots,b_n$ be real, with
$b=(b_1,\dots,b_n)\ne0$, and put $F(y)=\sum_{k=1}^nb_ke^{\lambda_ky}$. Then
$F$ has finitely many real zeros, and
$\operatorname{sp}(F(y_1),\dots,F(y_j))\preceq\operatorname{sp}(b)$ for all
$y_1<\dots<y_j$. The same holds for $\psi F$ for every positive function
$\psi$.
\end{lemma}

\begin{proof}
Let $Z(F)$ be the number of distinct real zeros of $F$, and $S(b)$ the number
of sign changes of $b$. We show $Z(F)\le S(b)$ by induction on $S(b)$. If
$S(b)=0$, then $F$ has the sign of the nonzero $b_k$ everywhere. If $S(b)\ge1$,
choose $i$ such that the last nonzero $b_k$ with $k\le i$ and the first
nonzero $b_k$ with $k>i$ have opposite signs, and choose
$\lambda_i<\kappa<\lambda_{i+1}$. Then
\[
 G(y)=\frac{d}{dy}\bigl(e^{-\kappa y}F(y)\bigr)
 =\sum_{k=1}^nb_k(\lambda_k-\kappa)e^{(\lambda_k-\kappa)y}.
\]
The coefficient $b_k(\lambda_k-\kappa)$ has the sign of $-b_k$ for $k\le i$
and of $b_k$ for $k>i$, so the coefficients of $G$ have $S(b)-1$ sign changes.
By Rolle's theorem $G$ vanishes between any two zeros of $F$, so
$Z(F)\le Z(G)+1\le S(b)$.

If $b_{k_0}$ is the first nonzero $b_k$, then
$e^{-\lambda_{k_0}y}F(y)\to b_{k_0}$ as $y\to-\infty$. The zeros of $F$ cut the
line into at most $S(b)+1$ open intervals, on each of which $F$ has one sign.
With equal neighbors merged, these signs form an alternating list that starts
with the sign of $b_{k_0}$ and has at most $S(b)+1$ entries, so it is
$\preceq\operatorname{sp}(b)$. The nonzero values $F(y_1),\dots,F(y_j)$ lie in
these intervals in order. A positive factor $\psi$ changes no sign.
\end{proof}

\begin{lemma}[pointwise limits]\label{lem:sp-limit}
Let $v_n$ and $v$ be real functions on $\mathbb R$ with $v_n(y)\to v(y)$ for
every $y$, and let $\Pi$ be an alternating list. If every $v_n$ has values at
most $\Pi$, then so has $v$. If every $v_n$ has increments at most $\Pi$, then
so has $v$.
\end{lemma}

\begin{proof}
Let $y_1<\dots<y_j$, and let $I$ be the set of indices $i$ with $v(y_i)\ne0$.
For all large $n$, $v_n(y_i)$ has the sign of $v(y_i)$ for every $i\in I$.
Then $\operatorname{sp}(v(y_1),\dots,v(y_j))$ is the pattern of the list
$(v_n(y_i))_{i\in I}$, which is a subsequence of
$\operatorname{sp}(v_n(y_1),\dots,v_n(y_j))\preceq\Pi$. The increments of $v_n$
at $y_1,\dots,y_j$ tend to those of $v$, and the same argument applies to
them.
\end{proof}

\begin{lemma}[one derivative fewer]\label{lem:sp-derivative}
Let $u$ be twice continuously differentiable on $\mathbb R$, with $u'(y)\to0$
as $y\to\pm\infty$. If $u''$ has values at most $\Pi_5$, then $u'$ has values
at most $\Pi_4$, and $u$ has increments at most $\Pi_4$.
\end{lemma}

\begin{proof}
It is enough to treat points $y_1<\dots<y_j$ at which the values of $u'$ are
nonzero and alternate in sign. Let $s$ be the sign of $u'(y_1)$. Choose $y_0<y_1$
with $\lvert u'(y_0)\rvert<\lvert u'(y_1)\rvert$, and $y_{j+1}>y_j$ with
$\lvert u'(y_{j+1})\rvert<\lvert u'(y_j)\rvert$. For $0\le i\le j$ the mean
value theorem gives $w_i\in(y_i,y_{i+1})$ with
$u''(w_i)(y_{i+1}-y_i)=u'(y_{i+1})-u'(y_i)$. The right side has the sign $s$
for $i=0$, the sign of $u'(y_{i+1})$ for $1\le i\le j-1$, and the sign of
$-u'(y_j)$ for $i=j$. In each case this is $(-1)^is$. So $u''$ takes $j+1$
alternating signs, starting with $s$, at $w_0<\dots<w_j$. This list is
$\preceq\Pi_5$, so $j+1\le5$, and $j+1\le4$ if $s$ is $-$. In both cases the
list of the $j$ signs of $u'(y_1),\dots,u'(y_j)$ is $\preceq\Pi_4$.

For the increments at $y_1<\dots<y_j$, each nonzero $u(y_{i+1})-u(y_i)$ has
the sign of $u'(\xi_i)$ for some $\xi_i\in(y_i,y_{i+1})$, by the mean value
theorem. So the pattern of the increments is a pattern of values of $u'$.
\end{proof}
